\section{Related Work}
\label{sec:related}
\subsection{Radio-Frequency Fingerprinting}
Early radiometric identification established that analogue transmitter imperfections can provide device-specific signatures~\cite{brik2008wireless_device_id}. Deep convolutional models and large experimental studies subsequently made RFFI practical across modulation and channel conditions~\cite{riyaz2018deep_cnn_radio_id,sankhe2019oracle,jian2020deeplearning_rf_fingerprinting}. LoRa-specific studies have explored spectrograms, coherent integration, data augmentation and channel-robust feature learning~\cite{shen2021lora_jsac,shen2022scalable_lora_rffi,xie2018coherent_rff_iotj,alshawabka2021deeplora}.

\subsection{OOB Evidence and Receiver Shift}
OOB spectral distortion is a physically plausible source of transmitter evidence. WideScan and related studies demonstrate that spectral regions outside the nominal signal band can improve device classification~\cite{elmaghbub2020widescan,hamdaoui2020deep_nn_features}. This evidence is not automatically receiver invariant. Deployment variability, receiver-agnostic training and source-free adaptation have therefore become active topics in wireless RFFI~\cite{elmaghbub2021lora_wild,shen2024receiveragnostic,ma2025receiver_independent_rffi,yang2025scrffi}. Existing work primarily optimizes transfer performance. Our focus is complementary: we measure how a model's OOB reliance changes under controlled interventions and whether the same dependency reproduces on held-out receivers.

\subsection{Position of This Study}
We make three distinctions that are often conflated. First, OOB utility is not equivalent to OOB robustness. Second, a controlled perturbation is a probe of a learned dependency, not a causal decomposition of all receiver shift. Third, packet count is not the independent unit for cross-receiver inference. We therefore report complete receiver-by-seed results, use strict source/blind separation, and treat the receiver as the top-level statistical unit.
